\textbf{Proof.}
By Theorem \(\mathrm{thm:signed\text{-}canonical\text{-}reduction}\), every model in Definition \(\mathrm{def:identification\text{-}criterion}\) can be replaced by a counterfactually equivalent canonical model satisfying \(|\Omega_u|\le b_u\) for all \(u\in U_G\).  This replacement preserves every hard-intervention menu law and both counterfactual probabilities entering \(Q\).  Pad a smaller support by zero-mass states.  For a fixed enumerated signed-isotone structural profile \(a\), write the exogenous product-simplex coordinates as
\[
 \lambda_{u,j}\ge0,
 \qquad
 \sum_{j=1}^{b_u}\lambda_{u,j}=1.
\tag{1}
\]
Eliminating \(\lambda_{u,b_u}\) leaves at most
\[
 N_{G,M}=\sum_{u\in U_G}(b_u-1)
\tag{2}
\]
free real coordinates for one model.

For \(\omega\in\prod_{u\in U_G}\{1,\ldots,b_u\}\), the probability of the joint exogenous state is the monomial
\[
 w_\lambda(\omega)=\prod_{u\in U_G}\lambda_{u,\omega_u}.
\tag{3}
\]
Fix a regime \(z_k=\operatorname{do}(X_{W_k}=w_k)\).  Replace the structural equation of every vertex in \(W_k\) by its declared constant in \(w_k\), retain every other equation from profile \(a\), and evaluate the remaining equations in a topological order.  Acyclicity makes the resulting vector \(X^{a,z_k}(\omega)\in\{0,1\}^{V}\) unique.  Consequently every full regime cell is the integer-coefficient polynomial
\[
 \phi_{a,k,x}(\lambda)
 =\sum_{\omega}
   \mathbf 1\{X^{a,z_k}(\omega)=x\}
   w_\lambda(\omega).
\tag{4}
\]
This is the step for which deterministic hard-intervention syntax is load-bearing.  No stochastic intervention kernel is being treated as an unstated rational coefficient array.

The same recursive evaluation determines every unnested literal in \(A_Q\) and \(B_Q\).  Hence the denominator and numerator are the integer-coefficient polynomials
\[
 d_a(\lambda)
 =\sum_{\omega}\mathbf 1\{B_Q(\omega)\}w_\lambda(\omega),
 \qquad
 h_a(\lambda)
 =\sum_{\omega}\mathbf 1\{A_Q(\omega)\cap B_Q(\omega)\}w_\lambda(\omega).
\tag{5}
\]
Every sum in (4)--(5) has at most
\(
 R_{G,M}=\prod_{u\in U_G}b_u
\)
monomials.

Let \(\Lambda_a(\lambda)\) denote the finite system (1).  Nonidentification is exactly the finite disjunction, over ordered profile pairs \((a_0,a_1)\), of
\[
 \exists\lambda^0,\lambda^1\;\left[
 \begin{aligned}
 &\Lambda_{a_0}(\lambda^0)\wedge\Lambda_{a_1}(\lambda^1),\\
 &d_{a_0}(\lambda^0)>0\wedge d_{a_1}(\lambda^1)>0,\\
 &\phi_{a_0,k,x}(\lambda^0)=\phi_{a_1,k,x}(\lambda^1)
       &&\text{for every \(k\) and \(x\)},\\
 &h_{a_0}(\lambda^0)d_{a_1}(\lambda^1)
   \ne h_{a_1}(\lambda^1)d_{a_0}(\lambda^0).
 \end{aligned}\right]
\tag{6}
\]
The last line is equivalent to unequal conditional query values because the preceding line makes both denominators strictly positive.  By (2), each disjunct has at most \(2N_{G,M}\) free mass coordinates.  Lemma \(\mathrm{lem:effective\text{-}real\text{-}closed\text{-}field\text{-}decision}\) decides (6) and extracts an exact algebraic sample point when it is satisfiable.

If (6) is satisfiable, the sampled profiles and product-simplex masses define \(\mathfrak m_0,\mathfrak m_1\in\mathfrak M(G,M)\).  The table enumeration enforces every signed local-monotonicity restriction, and (1) enforces the semi-Markovian product law.  Equations (4) and (6) give
\[
 p^star_{k,x}
 =\phi_{a_0,k,x}(\lambda^0)
 =\phi_{a_1,k,x}(\lambda^1)
\]
for every menu cell, while (5)--(6) give positive denominators and
\(
 \theta_Q(\mathfrak m_0)\ne\theta_Q(\mathfrak m_1)
\).
Each model uses at most \(R_{G,M}\) joint latent configurations, so the pair uses at most \(2R_{G,M}\).

Conversely, if \(\operatorname{GID}_{G,M,\mathcal Z}(Q)=0\), Definition \(\mathrm{def:identification\text{-}criterion}\) supplies two models with the same specified hard-intervention menu law, positive query denominators, and unequal query values.  Apply Theorem \(\mathrm{thm:signed\text{-}canonical\text{-}reduction}\) to each.  Preservation of all finite counterfactual laws preserves every menu cell in (4) and both quantities in (5), so the resulting canonical pair satisfies one disjunct of (6).  Thus (6) is satisfiable if and only if the global identification criterion is false.  Effective real quantifier elimination proves decidability, and exact sample-point extraction proves the asserted countermodel conclusion. \(\square\)